\section{이론적 배경}


\subsection{선행 연구}
페달링의 힘은 두 방향으로 나누어 분석된다. 하나는 페달 힘을 크랭크에 수직인 유효력과 나란한 성분으로 나누는 것이고, 다른 하나는 관절들이 일을 나누어 맡는 비율을 보는 것이다. 비원형 체인링 연구는 이 두 가지를 모두 쓴다(<Table~\ref{tab:prior}>). 본 연구는 2.3절의 힘 분해로 첫째 방법을 모델에 그대로 적용하고(4.4절), 둘째 결과로 모델이 맞는 케이던스 범위를 판단한다(5.1절).

\begin{table}[!htb]
\centering\small\renewcommand{\arraystretch}{1.05}
\begin{tabular}{@{}L{27mm} L{40mm} L{63mm}@{}}
\toprule
연구 & 분석 방법 & 결과\\\midrule
Candotti 등\cite{candotti2007} & 측정한 페달 힘을 크랭크 좌표로 돌려 유효력(크랭크에 수직)과 비유효력으로 분해, IE = 유효력/합력 & 사이클 선수의 IE가 철인3종 선수보다 높았고, 차이는 60, 75 rpm에서만 유의\\
Bini 등\cite{bini2013} & 위와 같은 분해 (리뷰) & IE는 출력이 높을수록 커지고 케이던스가 높을수록 작아짐. IE를 높이라고 지시하면 효율은 오히려 떨어짐\\
Aasvold 등\cite{aasvold2019} & 고관절, 무릎, 발목의 관절 일률 & 60 rpm을 넘으면 케이던스가 오를수록 무릎 몫이 늘고 고관절 몫이 줄어듦. 낮은 강도에서 60 rpm 이하는 고관절, 80 rpm 이상은 무릎이 주로 일함\\
Rankin, Neptune\cite{rankin2008} & 근골격 모형(다리당 근육 15개, 힘-길이-속도 관계, 활성화 지연)으로 크랭크 일률 최대화 & 최적 지름비 60 rpm 1.35, 90 rpm 1.29, 120 rpm 1.24, 일률 약 3\% 증가. 힘을 내는 아래쪽 행정에서 크랭크를 늦추는 것이 원리\\
Leong 등\cite{leong2017} & 측정한 페달 힘과 다리 움직임으로 관절별 일률 계산 & 지름비 1.13, 1.24에서 최대 출력과 관절별 일률 변화 없음. 발목 각속도를 바꿔 무릎과 고관절 움직임을 유지\\
\bottomrule
\end{tabular}
\caption{페달링 힘의 분해와 비원형 체인링에 관한 선행 연구}
\label{tab:prior}
\end{table}
